%==============================================================
% PARTE 1: PRIMEROS PASOS
%==============================================================
\chapter{Parte 1: Primeros Pasos}

%--------------------------------------------------------------
\section{Creación del espacio de trabajo}
%--------------------------------------------------------------

En esta sección se describe el proceso de creación y configuración del espacio de trabajo ROS 2 (\texttt{ros2\_ws}), así como la instalación del entorno Docker con ROS 2 Humble.

% Describe brevemente los pasos realizados y añade capturas si lo crees necesario.

%--------------------------------------------------------------
\section{Ejemplo de un entorno: Nodos y Topics}
%--------------------------------------------------------------

\subsection*{Pregunta 1}
\textit{¿Qué podemos saber de este nodo? (Suscriptores, topics)}

\textbf{Respuesta:}\\
% Escribe tu respuesta aquí.

% Si tienes captura, inclúyela así:
% \begin{figure}[H]
%   \centering
%   \includegraphics[width=0.85\textwidth]{parte1/imagenes/p1_nodo_info.png}
%   \caption{Información del nodo \texttt{/talker} con \texttt{ros2 node info}.}
%   \label{fig:p1_nodo_info}
% \end{figure}

\subsection*{Pregunta 2}
\textit{¿Qué podemos saber sobre el topic que publica el nodo que acabamos de lanzar?}

\textbf{Respuesta:}\\
% Escribe tu respuesta aquí.

% \begin{figure}[H]
%   \centering
%   \includegraphics[width=0.85\textwidth]{parte1/imagenes/p1_topic_info.png}
%   \caption{Información del topic \texttt{/chatter}.}
%   \label{fig:p1_topic_info}
% \end{figure}

\subsection*{Pregunta 3}
\textit{¿Qué podemos saber sobre el nuevo nodo \texttt{/listener}?}

\textbf{Respuesta:}\\
% Escribe tu respuesta aquí.

% \begin{figure}[H]
%   \centering
%   \includegraphics[width=0.85\textwidth]{parte1/imagenes/p1_listener_info.png}
%   \caption{Información del nodo \texttt{/listener}.}
%   \label{fig:p1_listener_info}
% \end{figure}

\subsection*{Pregunta 4}
\textit{¿Qué ha sucedido ahora en el topic \texttt{/chatter}?}

\textbf{Respuesta:}\\
% Escribe tu respuesta aquí.

% \begin{figure}[H]
%   \centering
%   \includegraphics[width=0.85\textwidth]{parte1/imagenes/p1_chatter_subs.png}
%   \caption{Estado del topic \texttt{/chatter} tras conectar el listener.}
%   \label{fig:p1_chatter_subs}
% \end{figure}

\subsection*{Pregunta 5}
\textit{¿Qué podemos saber a partir del gráfico de \texttt{rqt\_graph}?}

\textbf{Respuesta:}\\
% Escribe tu respuesta aquí.

\begin{figure}[H]
  \centering
  % \includegraphics[width=0.85\textwidth]{parte1/imagenes/p1_rqt_graph.png}
  \fbox{\parbox{0.8\textwidth}{\centering\vspace{2cm} [Captura rqt\_graph] \vspace{2cm}}}
  \caption{Gráfico \texttt{rqt\_graph} con los nodos \texttt{/talker} y \texttt{/listener}.}
  \label{fig:p1_rqt_graph}
\end{figure}

%--------------------------------------------------------------
\section{Servicios}
%--------------------------------------------------------------

\subsection*{Pregunta 6}
\textit{¿Qué puedes comentar acerca de este servicio?}

\textbf{Respuesta:}\\
% Escribe tu respuesta aquí.

% \begin{figure}[H]
%   \centering
%   \includegraphics[width=0.85\textwidth]{parte1/imagenes/p1_servicio_suma.png}
%   \caption{Ejecución del servicio de suma.}
%   \label{fig:p1_servicio_suma}
% \end{figure}

\subsection*{Pregunta 7}
\textit{Al lanzar el servidor y ejecutar \texttt{ros2 service list} ¿Qué servicios aparecen? ¿Cómo se ha creado el servicio \texttt{minimal\_service} y cuál es su cometido?}

\textbf{Respuesta:}\\
% Escribe tu respuesta aquí.

% \begin{figure}[H]
%   \centering
%   \includegraphics[width=0.85\textwidth]{parte1/imagenes/p1_service_list.png}
%   \caption{Listado de servicios con \texttt{ros2 service list}.}
%   \label{fig:p1_service_list}
% \end{figure}

\subsection*{Pregunta 8}
\textit{¿Cómo puedo saber cuáles son los parámetros de la ``petición'' y de la ``respuesta'' de un servicio? ¿De qué tipología son estos parámetros?}

\textbf{Respuesta:}\\
% Escribe tu respuesta aquí. Ejemplo de comando:
\begin{lstlisting}[style=bash, caption={Mostrar interfaz del servicio con ros2 interface show.}]
ros2 interface show <tipo_servicio>
\end{lstlisting}

% \begin{figure}[H]
%   \centering
%   \includegraphics[width=0.85\textwidth]{parte1/imagenes/p1_interface_show.png}
%   \caption{Salida de \texttt{ros2 interface show} para el servicio de suma.}
%   \label{fig:p1_interface_show}
% \end{figure}

%--------------------------------------------------------------
\section{Ejercicios}
%--------------------------------------------------------------

\subsection{Ejercicio 1: Publicar mensaje al topic \texttt{/chatter}}

\textit{Publica desde la terminal un mensaje al topic \texttt{/chatter} para que sea leído por el nodo \texttt{/listener}.}

\textbf{Comandos utilizados:}
\begin{lstlisting}[style=bash, caption={Publicación manual en el topic /chatter.}]
ros2 topic pub /chatter std_msgs/msg/String "data: 'Tu mensaje aquí'"
\end{lstlisting}

% \begin{figure}[H]
%   \centering
%   \includegraphics[width=0.85\textwidth]{parte1/imagenes/p1_ej1_pub.png}
%   \caption{Publicación de mensaje en \texttt{/chatter} y recepción en \texttt{/listener}.}
%   \label{fig:p1_ej1_pub}
% \end{figure}

\subsection{Ejercicio 2: Modificar el servicio suma}

\textit{Cambia el servicio \texttt{servicio\_suma} para que calcule: \(a + b - c \cdot d \div e\), siendo \(a, b, c, d, e\) valores enteros (\texttt{int64}) y el resultado en \texttt{float64}.}

\textbf{Descripción de los cambios:}\\
% Explica qué archivos has modificado y cómo.

\textbf{Archivos modificados:}
\begin{itemize}
  \item \texttt{interfaz\_servicio\_suma/srv/...} -- Modificación de la interfaz del servicio.
  \item \texttt{servicio\_suma/...} -- Modificación del servidor y cliente.
\end{itemize}

% \begin{figure}[H]
%   \centering
%   \includegraphics[width=0.85\textwidth]{parte1/imagenes/p1_ej2_resultado.png}
%   \caption{Resultado de la ejecución del servicio modificado.}
%   \label{fig:p1_ej2_resultado}
% \end{figure}
